\documentclass[14pt,a4paper,oneside]{extbook}
\usepackage[utf8]{vietnam}
\usepackage{amsmath,amssymb,mathrsfs,fancyhdr,enumerate,multirow,makecell}
\usepackage{tkz-euclide,tikz-3dplot,tikz,tkz-tab}
\usetikzlibrary{arrows.meta,calc,intersections,angles,quotes,backgrounds,shapes.geometric,patterns,through,shadows,positioning}
\usepackage{yhmath}
\usepackage[top=0.7cm, bottom=0.7cm, left=1.0cm, right=1.0cm, headheight=16pt, headsep=3mm, footskip=5mm]{geometry}
\usepackage[hidelinks,unicode]{hyperref}
\usepackage{qrcode}
\usepackage[many]{tcolorbox}
\usepackage{fontawesome}

\pagestyle{fancy}
\fancyhf{}
\fancyhead[L]{\footnotesize\bfseries\color{teal!85!black}\faBook~SỔ TAY KIẾN THỨC TOÁN 4}
\fancyhead[R]{\footnotesize\bfseries\color{blue!85!black}Lớp Toán Thầy Lĩnh -- loptoanthaylinh.vn}
\fancyfoot[C]{\footnotesize\bfseries\color{gray!80!black}-- Sổ tay kiến thức: Bài 59. Luyện tập chung phân số --}
\renewcommand{\headrulewidth}{0.4pt}

\begin{document}

% --- BANNER TIÊU ĐỀ BÀI HỌC ---
\noindent
\begin{tcolorbox}[
  enhanced,
  colback=teal!8!white,
  colframe=teal!80!black,
  arc=3mm, boxrule=1pt,
  left=3mm, right=3mm, top=2mm, bottom=2mm,
  halign=center
]
  {\Large\bfseries\color{teal!90!black}\faBookmark\ BÀI 59: LUYỆN TẬP CHUNG PHÂN SỐ}\\[1.5pt]
  {\footnotesize\itshape\color{black!75} Hệ thống hóa toàn bộ mạch kiến thức phân số, sắp xếp dãy số và giải toán thực tế trên 1 trang A4}
\end{tcolorbox}

\vspace{-2.5mm}

\noindent
\begin{minipage}[t]{0.49\textwidth}\vspace{0pt}
\begin{tcolorbox}[
  enhanced,
  colback=blue!4!white,
  colframe=blue!75!black,
  arc=2mm, boxrule=0.8pt,
  left=2mm, right=2mm, top=1.8mm, bottom=1.8mm,
  title={\footnotesize\bfseries\color{white}\faCube\ 1. Bốn kĩ năng trụ cột phân số}
]
{\footnotesize
Hệ thống hóa 4 kĩ năng không thể thiếu:
\begin{itemize}\setlength{\itemsep}{1pt}\setlength{\leftskip}{-2.5mm}
  \item \textbf{1. Đọc, viết \& ý nghĩa:} Tử số / Mẫu số ($b > 0$).
  \item \textbf{2. Phân số bằng nhau:} Nhân/chia cả tử và mẫu với cùng số $m \ne 0$.
  \item \textbf{3. Rút gọn:} Chia cả tử và mẫu cho ước chung để về \textbf{phân số tối giản}.
  \item \textbf{4. Quy đồng \& So sánh:} Tìm MSC tối ưu để so sánh chính xác.
\end{itemize}
}
\end{tcolorbox}
\end{minipage}
\hfill
\begin{minipage}[t]{0.49\textwidth}\vspace{0pt}
\begin{tcolorbox}[
  enhanced,
  colback=teal!4!white,
  colframe=teal!75!black,
  arc=2mm, boxrule=0.8pt,
  left=2mm, right=2mm, top=1.8mm, bottom=1.8mm,
  title={\footnotesize\bfseries\color{white}\faSortAmountAsc\ 2. Chiến lược sắp xếp dãy phân số}
]
{\footnotesize
Để sắp xếp dãy nhiều phân số từ bé đến lớn:
\begin{itemize}\setlength{\itemsep}{1pt}\setlength{\leftskip}{-2.5mm}
  \item \textbf{Bước 1 (Phân loại theo 1):} Tách thành 2 nhóm: Nhóm bé hơn 1 và nhóm lớn hơn 1.
  \item \textbf{Bước 2 (Tối giản):} Rút gọn các phân số chưa tối giản để số nhỏ gọn.
  \item \textbf{Bước 3 (Quy đồng theo nhóm):} Quy đồng mẫu các phân số trong từng nhóm để so sánh.
\end{itemize}
}
\end{tcolorbox}
\end{minipage}

\vspace{-2.5mm}
% --- KHỐI 3: SƠ ĐỒ TƯ DUY TIKZ ---
\noindent
\begin{tcolorbox}[
  enhanced,
  colback=orange!4!white,
  colframe=orange!80!black,
  arc=2mm, boxrule=0.8pt,
  left=2mm, right=2mm, top=1.5mm, bottom=1.5mm,
  title={\footnotesize\bfseries\color{white}\faSitemap\ SƠ ĐỒ TƯ DUY: BẢN ĐỒ MẠNG NHỆN 4 CHÌA KHÓA PHÂN SỐ}
]
\centering
\begin{tikzpicture}[
  box/.style={draw, rounded corners=2pt, align=center, font=\scriptsize\bfseries, inner sep=3pt, minimum width=3.4cm, minimum height=0.9cm},
  arr/.style={-{Stealth[length=2mm]}, thick}
]
  \node[draw=purple!80!black, fill=purple!15, rounded corners=3pt, font=\small\bfseries\color{purple!90!black}, inner sep=4pt, align=center] (core) at (5, 0.7) {MẠCH KIẾN THỨC\\PHÂN SỐ};

  \node[box, fill=blue!10, draw=blue!80!black] (k1) at (0.8, 1.4) {1. KHÁI NIỆM\\Tử / Mẫu ($b \ne 0$)};
  \node[box, fill=teal!10, draw=teal!80!black] (k2) at (9.2, 1.4) {2. TÍNH CHẤT\\Nhân/chia cùng số};
  \node[box, fill=orange!12, draw=orange!80!black] (k3) at (0.8, 0.0) {3. RÚT GỌN\\Đưa về Tối giản};
  \node[box, fill=green!15, draw=green!70!black] (k4) at (9.2, 0.0) {4. QUY ĐỒNG\\So sánh \& Sắp xếp};

  \draw[arr, blue!80!black] (core) -- (k1);
  \draw[arr, teal!80!black] (core) -- (k2);
  \draw[arr, orange!80!black] (core) -- (k3);
  \draw[arr, green!70!black] (core) -- (k4);
\end{tikzpicture}
\end{tcolorbox}

\vspace{-2.5mm}
% --- KHỐI 4: VÍ DỤ & MẸO TRÁNH BẪY ---
\noindent
\begin{minipage}[t]{0.49\textwidth}\vspace{0pt}
\begin{tcolorbox}[
  enhanced,
  colback=green!4!white,
  colframe=green!60!black,
  arc=2mm, boxrule=0.8pt,
  left=2mm, right=2mm, top=1.8mm, bottom=1.8mm,
  title={\footnotesize\bfseries\color{white}\faCheckCircle\ Ví dụ thực hành chuẩn mực}
]
{\footnotesize
\textbf{Ví dụ: Sắp xếp các phân số sau từ bé đến lớn:}\\
$$\dfrac{1}{2}; \quad \dfrac{3}{4}; \quad \dfrac{5}{4}; \quad \dfrac{1}{4}$$
\begin{itemize}\setlength{\itemsep}{1pt}\setlength{\leftskip}{-2.5mm}
  \item Phân loại: $\dfrac{5}{4} > 1$ (lớn nhất).
  \item Ba phân số còn lại cùng $< 1$. Đưa về cùng mẫu số 4: $\dfrac{1}{2} = \dfrac{2}{4}$.
  \item So sánh: $\dfrac{1}{4} < \dfrac{2}{4} < \dfrac{3}{4} < \dfrac{5}{4}$.
  \item Kết luận: $\mathbf{\dfrac{1}{4} < \dfrac{1}{2} < \dfrac{3}{4} < \dfrac{5}{4}}$.
\end{itemize}
}
\end{tcolorbox}
\end{minipage}
\hfill
\begin{minipage}[t]{0.49\textwidth}\vspace{0pt}
\begin{tcolorbox}[
  enhanced,
  colback=red!4!white,
  colframe=red!70!black,
  arc=2mm, boxrule=0.8pt,
  left=2mm, right=2mm, top=1.8mm, bottom=1.8mm,
  title={\footnotesize\bfseries\color{white}\faExclamationTriangle\ Mẹo tránh bẫy trắc nghiệm}
]
{\footnotesize
\begin{itemize}\setlength{\itemsep}{1.5pt}\setlength{\leftskip}{-2.5mm}
  \item \textbf{Bẫy quy đồng tất cả các mẫu:}\\
  Nếu dãy có 4 phân số mà nhân hết 4 mẫu lại thì MSC cực kì lớn. Luôn tìm mẫu số chung nhỏ nhất hoặc phân nhóm so sánh qua 1.
  \item \textbf{Bẫy quên viết lại phân số gốc:}\\
  Sau khi quy đồng để so sánh, khi trả lời kết luận phải dùng các \textbf{phân số ban đầu} của đề bài.
\end{itemize}
}
\end{tcolorbox}
\end{minipage}

\end{document}
